\section{AIME I 1984}
        \begin{enumerate}
        	\item (\textit{AIME I 1984}) Find the value of $a_2+a_4+a_6+a_8+\ldots+a_{98}$ if $a_1$, $a_2$, $a_3\ldots$ is an arithmetic progression with common difference 1, and $a_1+a_2+a_3+\ldots+a_{98}=137$.


 
 \textbf{Answer: }093



\textbf{Solution 1}

One approach is to sum an \href{https://artofproblemsolving.com/wiki/index.php?title=Arithmetic\_series}{arithmetic series}, find the value of $a_1$, calculate $a_2$, and sum another arithmetic series to get the answer. A quicker method is to substitue $a_{2n - 1} = a_{2n} - 1$ into the given equation: 
 \[(a_2 - 1) + a_2 + (a_4 - 1) + a_4 + \ldots + (a_{98} - 1) + a_{98} = 137\]Simplify the LHS:
 \[2(a_2 + a_4 + \ldots + a_{98}) - 49= 137\]Solve:
 \[a_2 + a_4 + \ldots + a_{98}=\frac{137 + 49}{2} = \boxed{93}\]

 



\textbf{Solution 2}

If $a_1$ is the first term, then $a_1+a_2+a_3 + \cdots + a_{98} = 137$ can be rewritten as:

 $98a_1 + 1+2+3+ \cdots + 97 = 137$ $\Leftrightarrow$$98a_1 + \frac{97 \cdot 98}{2} = 137$

 Our desired value is $a_2+a_4+a_6+ \cdots + a_{98}$ so this is:

 $49a_1 + 1+3+5+ \cdots + 97$

 which is $49a_1+ 49^2$. So, from the first equation, we know $49a_1 = \frac{137}{2} - \frac{97 \cdot 49}{2}$. So, the final answer is:

 $\frac{137 - 97(49) + 2(49)^2}{2} = \fbox{93}$.

 



\textbf{Solution 3}

A better approach to this problem is to notice that from $a_{1}+a_{2}+\cdots a_{98}=137$ that each element with an odd subscript is 1 from each element with an even subscript. Thus, we note that the sum of the odd elements must be $\frac{137-49}{2}$. Thus, if we want to find the sum of all of the even elements we simply add $49$ common differences to this giving us $\frac{137-49}{2}+49=\fbox{93}$.

 Or, since the sum of the odd elements is 44, then the sum of the even terms must be $\fbox{93}$.

 



\textbf{Solution 4}

We want to find the value of $a_2+a_4+a_6+a_8+\ldots+a_{98}$, which can be rewritten as $a_1+1+a_2+2+a_3+\ldots+a_{49}+49 \implies a_1+a_2+a_3+\ldots+a_{49}+\frac{49 \cdot 50}{2}$.We can split $a_1+a_2+a_3+\ldots+a_{98}$ into two parts:
 \[a_1+a_2+a_3+\ldots+a_{49}\] and 
 \[a_{50}+a_{51}+a_{52}+\ldots+a_{98}\]Note that each term in the second expression is $49$ greater than the corresponding term, so, letting the first equation be equal to $x$, we get $a_1+a_2+a_3+\ldots+a_{98}=137=2x+49^2 \implies x=\frac{137-49^2}{2}$. Calculating $49^2$ by sheer multiplication is not difficult, but you can also do $(50-1)(50-1)=2500-100+1=2401$. We want to find the value of $x+\frac{49 \cdot 50}{2}=x+49 \cdot 25=x+1225$. Since $x=\frac{137-2401}{2}$, we find $x=-1132$. $-1132+1225=\boxed{93}$.

 - PhunsukhWangdu

 



\textbf{Solution 5}

Since we are dealing with an arithmetic sequence, 
 \[a_2+a_4+a_6+a_8+\ldots+a_{98} = 49a_{50}\]We can also figure out that 
 \[a_1+a_2+a_3+\ldots+a_{98} = a_1 + 97a_{50} = 137\]
 \[a_1 = a_{50}-49 \Rightarrow 98a_{50}-49 = 137\]Thus, $49a_{50} = \frac{137 + 49}{2} = \boxed{93}$

 



\textbf{Solution 6}

The sum of an arithmetic sequence of $n$ terms is $S_n = \frac{n}{2}(a_1 + a_n)$.

 We are told that the sum of all 98 terms is 137, so$\frac{98}{2}(a_1 + a_n) = 137$

 The subset of terms $a_2, a_4, a_6, a_8, ...$ themselves also form an arithmetic sequence (having common difference 2), so their sum $S$ will be$\frac{49}{2}(a_2 + a_{98}) = S$

 But $a_2$ is just $a_1 + 1$, so$\frac{49}{2}(a_1 + 1 + a_{98}) = S$

 Here we have a system of two equations. Manipulating the first equation yields$a_1 + a_{98} = \frac{137}{49}$

 And putting this into the second equation,$S = \frac{49}{2}(1 + \frac{137}{49}) = \fbox{093}$

 


	\item (\textit{AIME I 1984}) The integer $n$ is the smallest positive multiple of $15$ such that every digit of $n$ is either $8$ or $0$. Compute $\frac{n}{15}$.


 
 \textbf{Answer: }592



\textbf{Solution 1}

Any multiple of 15 is a multiple of 5 and a multiple of 3. 

 Any multiple of 5 ends in 0 or 5; since $n$ only contains the digits 0 and 8, the units \href{https://artofproblemsolving.com/wiki/index.php?title=Digit}{digit} of $n$ must be 0.  

 


 
The sum of the digits of any multiple of 3 must be \href{https://artofproblemsolving.com/wiki/index.php?title=Divisible}{divisible} by 3.  If $n$ has $a$ digits equal to 8, the sum of the digits of $n$ is $8a$.  For this number to be divisible by 3, $a$ must be divisible by 3. We also know that $a>0$ since $n$ is positive. Thus $n$ must have at least three copies of the digit 8.  

 The smallest number which meets these two requirements is 8880.  Thus the answer is $\frac{8880}{15} = \boxed{592}$.

 



\textbf{Solution 2}

Notice how $8 \cdot 10^k \equiv 8 \cdot (-5)^k \equiv 5 \pmod{15}$ for all integers $k \geq 2$. Since we are restricted to only the digits $8,0$, because $8\equiv -7 \pmod{15}$ we can't have an $8$ in the optimal smallest number. We can just 'add' fives to quickly get $15 \equiv 0 \pmod{15}$ to get our answer. Thus n is $80+800+8000=8880$ and $n/15=8880/15=\boxed{592}$

 



\textbf{Solution 3}

Any number that consists only of the digits $8$ and $0$ is itself divisible by $8$. $n$ therefore is divisible by $8$. Furthermore, $n$ must end in $0$ since it is divisible by $15$, and all multiples of $15$ will end in $5$ or $0$. Therefore may write $2^3 \cdot 3 \cdot 5 \cdot k = 15m$ to express that $n$ is simultaneously a multiple of $8$ and $15$. Going up by multiples of $120$, among three digit numbers we only have to check $840$ which clearly does not satisfy the condition of our problem. Among four digit numbers, you encounter $8040, 8160, 8280, 8400, 8520, 8640, 8760, 8880$ -- this last number corresponds to $k = 74$ or $m =  \boxed{592}$

 


	\item (\textit{AIME I 1984}) A point $P$ is chosen in the interior of $\triangle ABC$ such that when lines are drawn through $P$ parallel to the sides of $\triangle ABC$, the resulting smaller triangles $t_{1}$, $t_{2}$, and $t_{3}$ in the figure, have areas $4$, $9$, and $49$, respectively. Find the area of $\triangle ABC$.


 
                    \begin{figure}[H]
                        \centering
                        \includegraphics[width=0.5\textwidth]{Problem/AIME_I_1984_3_Problem_Figure_1.png}
                    \end{figure}
                    


 
 \textbf{Answer: }144



\textbf{Solution 1}

By the transversals that go through $P$, all four triangles are \href{https://artofproblemsolving.com/wiki/index.php?title=Similar\_triangles}{similar} to each other by the $AA$ postulate. Also, note that the length of any one side of the larger triangle is equal to the sum of the sides of each of the corresponding sides on the smaller triangles. We use the identity $K = \dfrac{ab\sin C}{2}$ to show that the areas are proportional (the sides are proportional and the angles are equal)

 



\textbf{Solution 2}

Alternatively, since the triangles are similar by $AA$, then the ratios between the bases and the heights of each of the three triangles would all be equal. The areas of each of the triangles are all perfect squares, so we could assume $\dfrac{base}{height} = \dfrac{2}{1}.$ That means that the base of $t_{1}$ is 4, the base of $t_{2}$ is 6, and the base of $t_{3}$ is 14. Since the quadrilaterals underneath $t_{1}$ and $t_{2}$ are both parallelograms, and opposite sides of a parallelogram are congruent, the base of the large triangle is $4 + 14 + 6 = 24$. Therefore, the height of the entire triangle would be twelve, so therefore, the area of the large triangle is $\dfrac{1}{2} \cdot 24 \cdot 12 = \boxed{144}$.

 



\textbf{Solution 3}

The base of $\triangle{ABC}$ is $BC$. Let the base of $t_1$ be $x$, the base of $t_2$ be $y$, and the base of $t_3$ be $z$. Since $\triangle{ABC}, t_1, t_2,$ and $t_3$ are all similar, the sections in $\triangle{ABC}$ that aren't $t_1,t_2,$ or $t_3$ are all parallelograms. Hence, $BC=x+z+y$. We can relate $t_1,t_2,$ and $t_3$ by the square root of the ratio of their areas. $\sqrt{\frac{4}{9}}=\frac{2}{3}$ and $\sqrt{\frac{4}{49}}=\frac{2}{7}$ so $y=\frac{3x}{2}$ and $z=\frac{7x}{2}$. $x+\frac{7x}{2}+\frac{3x}{2}=6x$, so $\triangle{ABC}$ has a base that is $6$ times $t_1$. $[\triangle{ABC}]=36[t_1]=36 \cdot 4=\boxed{144}$.

 -PhunsukhWangdu

 


 



\textbf{Solution 4}

Since the three lines through $P$ are parallel to the sides, $t_1$, $t_2$, $t_3$, and $\triangle{ABC}$ are similar by $AA$ similarity. Suppose the area of $\triangle{ABC}$ is $x^2$, so the ratio of the base of $t_1$ to the base of $t_2$ to the base of $t_3$ to the base of $\triangle{ABC}$ is $2:3:7:x$. Because the quadrilaterals below $t_1$ and $t_2$ are parallelograms, the base of $\triangle{ABC}$ is equal to the sum of the bases of $t_1, t_2,$ and $t_3$. Therefore, $x$ equals $2+3+7=12$ so the area of $\triangle{ABC}$ equals $x^2=12^2=\boxed{144}.$

 -Yiyj1

 



\textbf{Solution 5}

Let the longest leg of $t3$ be $x$, and the other leg be $\frac{98}{x}$; we see that all $3$ triangles are similar, and we take the square root of their area ratios to find that the side ratio of $t3$ to $t2$ is $\frac{7}{3}$ and the ratio of $t3$ to $t1$ is $\frac{7}{2}$.

 We see that the two legs of $t2$ will be $\frac{3x}{7}$ and $\frac{42}{x}$, the two legs of $t1$ will be $\frac{2x}{7}$ and $\frac{28}{x}$. The rectangle above these two triangles will have an area of $\frac{2x}{7}$ times  $\frac{42}{x}$, or $12$. We can now create two rectangles, each with a base being one leg of $t3$ and one with a height of the longer leg of $t1$, and the other with a height of the shorter leg of $t2$. To create these rectangles, we get rid of 2 triangles from the bottom left and right of the figure. We assume these triangles are congruent to $t2$ and $t1$, respectively.

 The area of rectangle one, with a height of the longest leg of $t1$ and a base of the shortest leg of $t3$, will be $\frac{98}{x}$ times $\frac{2x}{7}$, or $28$. 

 The area of the second rectangle will be $x$ times $\frac{42}{x}$, or $42$. The remaining $2$ triangles will have a combined area of $13$ as they are identical to $t1$ and $t2$. We now sum the areas of the top rectangle, our $2$ rectangles, $t3$, and the bottom $2$ triangles.

 $12+28+42+49+13=\boxed{144}$

 


	\item (\textit{AIME I 1984}) Let $S$ be a list of positive integers--not necessarily distinct--in which the number $68$ appears. The average (arithmetic mean) of the numbers in $S$ is $56$. However, if $68$ is removed, the average of the remaining numbers drops to $55$. What is the largest number that can appear in $S$?


 
 \textbf{Answer: }649



\textbf{Solution 1 (Two Variables)}

Suppose that $S$ has $n$ numbers other than the $68,$ and the sum of these numbers is $s.$

 We are given that 
 \begin{align*} \frac{s+68}{n+1}&=56, \\ \frac{s}{n}&=55. \end{align*}
Clearing denominators, we have
 \begin{align*} s+68&=56n+56, \\ s&=55n. \end{align*}
Subtracting the equations, we get $68=n+56,$ from which $n=12.$ It follows that $s=660.$

 The sum of the twelve remaining numbers in $S$ is $660.$ To maximize the largest number, we minimize the other eleven numbers: We can have eleven $1$s and one $660-11\cdot1=\boxed{649}.$

 



\textbf{Solution 2 (One Variable)}

Suppose that $S$ has $n$ numbers other than the $68.$ We have the following table:
 \[\begin{array}{c|c|c|c} & & & \\ [-2.5ex] & \textbf{Count} & \textbf{Arithmetic Mean} & \textbf{Sum} \\ \hline & & & \\ [-2.5ex] \textbf{Initial} & n+1 & 56 & 56(n+1) \\ \hline & & & \\ [-2.5ex] \textbf{Final} & n & 55 & 55n \end{array}\]We are given that 
 \[56(n+1)-68=55n,\] from which $n=12.$ It follows that the sum of the remaining numbers in $S$ is $55n=660.$ We continue with the last paragraph of Solution 1 to get the answer $\boxed{649}.$

 


	\item (\textit{AIME I 1984}) Determine the value of $ab$ if $\log_8a+\log_4b^2=5$ and $\log_8b+\log_4a^2=7$.


 
 \textbf{Answer: }512



\textbf{Solution 1}

Use the \href{https://artofproblemsolving.com/wiki/index.php?title=Change\_of\_base\_formula}{change of base formula} to see that $\frac{\log a}{\log 8} + \frac{2 \log b}{\log 4} = 5$; combine \href{https://artofproblemsolving.com/wiki/index.php?title=Denominator}{denominators} to find that $\frac{\log ab^3}{3\log 2} = 5$. Doing the same thing with the second equation yields that $\frac{\log a^3b}{3\log 2} = 7$. This means that $\log ab^3 = 15\log 2 \Longrightarrow ab^3 = 2^{15}$ and that $\log a^3 b = 21\log 2 \Longrightarrow a^3 b = 2^{21}$. If we multiply the two equations together, we get that $a^4b^4 = 2^{36}$, so taking the fourth root of that, $ab = 2^9 = \boxed{512}$.

 



\textbf{Solution 2}

We can simplify our expressions by changing everything to a common base and by pulling exponents out of the logarithms. The given equations then become $\frac{\ln a}{\ln 8} + \frac{2 \ln b}{\ln 4} = 5$ and $\frac{\ln b}{\ln 8} + \frac{2 \ln a}{\ln 4} = 7$. Adding the equations and factoring, we get $(\frac{1}{\ln 8}+\frac{2}{\ln 4})(\ln a+ \ln b)=12$. Rearranging we see that $\ln ab = \frac{12}{\frac{1}{\ln 8}+\frac{2}{\ln 4}}$. Again, we pull exponents out of our logarithms to get $\ln ab = \frac{12}{\frac{1}{3 \ln 2} + \frac{2}{2 \ln 2}} = \frac{12 \ln 2}{\frac{1}{3} + 1} = \frac{12 \ln 2}{\frac{4}{3}} = 9 \ln 2$. This means that $\frac{\ln ab}{\ln 2} = 9$. The left-hand side can be interpreted as a base-2 logarithm, giving us $ab = 2^9 = \boxed{512}$.

 



\textbf{Solution 3}

This solution is very similar to the above two, but it utilizes the well-known fact that $\log_{m^k}{n^k}= \log_m{n}.$ Thus, $\log_8a+\log_4b^2=5 \Rightarrow  \log_{2^3}{(\sqrt[3]{a})^3} + \log_{2^2}{b^2} = 5 \Rightarrow \log_2{\sqrt[3]{a}} + \log_2{b} = 5 \Rightarrow \log_2{\sqrt[3]{a}b} = 5.$ Similarly, $\log_8b+\log_4a^2=7 \Rightarrow \log_2{\sqrt[3]{b}a} = 7.$ Adding these two equations, we have $\log_2{a^{\frac{4}{3}}b^{\frac{4}{3}}} = 12 \Rightarrow ab = 2^{12\times\frac{3}{4}} = 2^9 = \boxed{512}$.

 



\textbf{Solution 4}

We can change everything to a common base, like so: $\log_8{a} + \log_8{b^3} = 5,$ $\log_8{b} + \log_8{a^3} = 7.$ We set the value of $\log_8{a}$ to $x$, and the value of $\log_8{b}$ to $y.$ Now we have a system of linear equations: 
 \[x + 3y = 5,\] 
 \[y + 3x = 7.\] Now add the two equations together then simplify, we'll get $x+y=3$. So $\log_8{ab} = \log_8{a} + \log_8{b} = 3$, $ab = 8^3 = \boxed{512}$

 



\textbf{Solution 5}

Add the two equations to get $\log_8 {a}+ \log_8 {b}+ \log_{a^2}+\log_{b^2}=12$. This can be simplified with the log property $\log_n {x}+\log_n {y}=log_n {xy}$. Using this, we get $\log_8 {ab}+ \log_4 {a^2b^2}=12$. Now let $\log_8 {ab}=c$ and $\log_4 {a^2b^2}=k$. Converting to exponents, we get $8^c=ab$ and $4^k=(ab)^2$. Sub in the $8^c$ to get $k=3c$. So now we have that $k+c=12$ and $k=3c$ which gives $c=3$, $k=9$. This means $\log_4 {a^2b^2}=9$ so $4^9=(ab)^2 \implies ab=(2^2)^9 \implies 2^9 \implies \boxed {512}$

 



\textbf{Solution 6}

Add the equations and use the facts that $\log{a} +\log{b}=\log{ab}$ and $\log{k^n} =n\log{k}$ to get
 \[\log_8{ab} +2\log_4{ab}=12\]Now use the change of base identity with base as 2:
 \[\dfrac{\log_2{ab}}{\log_2{8}}+\dfrac{2\log_2{ab}}{\log_2{4}}=12\]Which gives:
 \[\frac{4}{3}\log_2{ab}=12\]Solving gives, $\boxed{ab=2^9=512}$

 



\textbf{Solution 7}

By properties of logarithms, we know that $\log_8 {a}+ \log_4 {b ^ 2} = \log_2 {a ^ {1/3}}+ \log_2 {b} = 5$. 

 Using the fact that $\log_a {b}  + \log_a {c}  = log_a{b*c}$, we get $\log_2 {a^{1/3} * b} = 5$. 

 Similarly, we know that $\log_2 {a * b^{1/3}} = 7$. 

 From these two equations, we get $a^{1/3} * b = 2^5$ and $a * b^{1/3} = 2^7$. 

 Multiply the two equations to get $a^{4/3} * b^{4/3} = 2^{12}$. Solving, we get that $a*b = 2^{12*3/4} = 2^9 =$$\boxed{512}$.

 



\textbf{Solution 8}

Adding both of the equations, we get 
 \[\log_8{ab} +2\log_4{ab}=12\]Furthermore, we see that $\log_4 {ab}$ is $\frac{3}{2}$ times $\log_8 {ab}.$ Substituting $\log_8 {ab}$ as $x,$ we get $x+3x=12,$ so $x=3.$ Therefore, we have $\log_8 {ab} = 3,$ so $ab= 8^3=\boxed{512}$ ~ math\_comb01

 



\textbf{Solution 9}

Change all equations to base 64. We then get:
 \[\log_{64}(a^2) + \log_{64}(b^6) = 5\]and
 \[\log_{64}(b^2) + \log_{64}(a^6) = 7.\]

 Using the property \(\log(a) + \log(b) = \log(ab)\), we get:
 \[\log_{64}(a^2b^6) = 5\]and
 \[\log_{64}(a^6b^2) = 7.\]

 Then:
 \[a^2b^6 = 64^5\]and
 \[a^6b^2 = 64^7.\]

 Simplifying, we have:
 \[ab^3 = 8^5\]and
 \[a^3b = 8^7.\]

 Substituting and solving, we get:
 \[a = 8^2\]and
 \[b = 8.\]

 Then:
 \[ab = 8^3 = 512.\]

 



\textbf{Solution 10}

Given:- $\log_8 a + \log_4 b^2 = 5$- $\log_8 b + \log_4 a^2 = 7$

 We set up a system by subtracting 2 from the second equation to make both equal to 5:- $\log_8 a + \log_4 b^2 = 5$- $\log_8 b + \log_4 a^2 - 2 = 5$

 Setting these equal:$\log_8 a + \log_4 b^2 = \log_8 b + \log_4 a^2 - 2$

 Using the property $\log_4 x^2 = 2\log_4 x$:$\log_8 a + 2\log_4 b = \log_8 b + 2\log_4 a - 2$

 Converting to base 2, where $\log_8 x = \frac{\log_2 x}{3}$ and $\log_4 x = \frac{\log_2 x}{2}$:$\frac{\log_2 a}{3} + 2\cdot\frac{\log_2 b}{2} = \frac{\log_2 b}{3} + 2\cdot\frac{\log_2 a}{2} - 2$

 Simplifying:$\frac{\log_2 a}{3} + \log_2 b = \frac{\log_2 b}{3} + \log_2 a - 2$

 Multiplying through by 6 to eliminate fractions:$2\log_2 a + 6\log_2 b = 2\log_2 b + 6\log_2 a - 12$

 Rearranging:$2\log_2 a - 6\log_2 a + 6\log_2 b - 2\log_2 b = -12$$-4\log_2 a + 4\log_2 b = -12$$\log_2 b - \log_2 a = -3$

 Therefore:$\log_2\left(\frac{b}{a}\right) = -3$$\frac{b}{a} = \frac{1}{8}$$b = \frac{a}{8}$

 From the first original equation:$\log_8 a + \log_4 b^2 = 5$

 Substituting $b = \frac{a}{8}$:$\log_8 a + \log_4 \left(\frac{a}{8}\right)^2 = 5$$\log_8 a + \log_4 \left(\frac{a^2}{64}\right) = 5$$\log_8 a + \log_4 a^2 - \log_4 64 = 5$

 Since $\log_4 64 = \log_4 4^3 = 3$:$\log_8 a + \log_4 a^2 - 3 = 5$$\log_8 a + \log_4 a^2 = 8$

 Now we have:- $\log_8 a + \log_4 a^2 = 8$ (derived)- $\log_8 b + \log_4 a^2 = 7$ (given)

 Subtracting:$\log_8 a - \log_8 b = 1$$\log_8\left(\frac{a}{b}\right) = 1$$\frac{a}{b} = 8$

 This confirms our earlier finding that $b = \frac{a}{8}$, so $a = 8b$.

 Substituting this back into the first original equation:$\log_8 (8b) + \log_4 b^2 = 5$$\log_8 8 + \log_8 b + \log_4 b^2 = 5$$1 + \log_8 b + \log_4 b^2 = 5$$\log_8 b + \log_4 b^2 = 4$

 From the second original equation:$\log_8 b + \log_4 (8b)^2 = 7$$\log_8 b + \log_4 (64b^2) = 7$$\log_8 b + \log_4 64 + \log_4 b^2 = 7$$\log_8 b + 3 + \log_4 b^2 = 7$$\log_8 b + \log_4 b^2 = 4$

 Thus both equations yield the same constraint: $\log_8 b + \log_4 b^2 = 4$

 Converting to base 2:$\frac{\log_2 b}{3} + \frac{2\log_2 b}{2} = 4$$\frac{\log_2 b}{3} + \log_2 b = 4$$\frac{\log_2 b + 3\log_2 b}{3} = 4$$\frac{4\log_2 b}{3} = 4$$\log_2 b = 3$

 Therefore:$b = 2^3 = 8$$a = 8b = 8 \cdot 8 = 64$$ab = 64 \cdot 8 = 512$

 The value of $ab$ is $\boxed{512}$.

 



\textbf{Solution 11}

Let us convert everything from the first equation to log base-2. We get $\frac{\log_2{a}}{3} + \frac{\log_2{b^2}}{2} = 5$. Multiplying by $6$, we get $2\log_2{a} + 3\log_2{b^2} = 30$. By the properties of logarithms, $\log_2{a^2b^6} = 30$, so $a^2b^6 = 2^{30}$. We can use the exact same process for the next equation. We have $\log_2{b^2a^6} = 42$, so $a^6b^2 = 2^{42}$. We can now multiply our two equations of products of powers of $a$ and $b$ to get $a^8b^8 = 2^{42} \cdot 2^{30} = 2^{72}$. Taking the $8$-th root of both sides, we get $ab = 2^{\frac{72}{8}} = 2^9 = \boxed{512}$

 


	\item (\textit{AIME I 1984}) Three circles, each of radius $3$, are drawn with centers at $(14, 92)$, $(17, 76)$, and $(19, 84)$. A line passing through $(17,76)$ is such that the total area of the parts of the three circles to one side of the line is equal to the total area of the parts of the three circles to the other side of it. What is the absolute value of the slope of this line?


 
 \textbf{Answer: }024



\textbf{Solution 1}


                    \begin{figure}[H]
                        \centering
                        \includegraphics[width=0.5\textwidth]{Solution/AIME_I_1984_6_Solution_Figure_1.png}
                    \end{figure}
                    

 The line passes through the center of the bottom circle; hence it is the circle's \href{https://artofproblemsolving.com/wiki/index.php?title=Diameter}{diameter} and splits the circle into two equal areas. For the rest of the problem, we do not have to worry about that circle. 

 Draw the \href{https://artofproblemsolving.com/wiki/index.php?title=Midpoint}{midpoint} of $\overline{AC}$ (the centers of the other two circles), and call it $M$. If we draw the feet of the \href{https://artofproblemsolving.com/wiki/index.php?title=Perpendicular}{perpendiculars} from $A,C$ to the line (call $E,F$), we see that $\triangle AEM\cong \triangle CFM$ by ASA congruency; hence $M$ lies on the line. The coordinates of $M$ are $\left(\frac{19+14}{2},\frac{84+92}{2}\right) = \left(\frac{33}{2},88\right)$. 

 Thus, the slope of the line is $\frac{88 - 76}{\frac{33}{2} - 17} = -24$, and the answer is $\boxed{024}$.

 \textit{Remark}: Notice the fact that the radius is 3 is not used in this problem; in fact changing the radius does not affect the answer.

 Minor edits by \href{https://artofproblemsolving.com/wiki/index.php?title=User:T3chn0b14d3&action=edit&redlink=1}{T3CHN0B14D3}

 



\textbf{Solution 2}

First of all, we can translate everything downwards by $76$ and to the left by $14$. Then, note that a line passing through a given point intersecting a circle with a center as that given point will always cut the circle in half, so we can re-phrase the problem:

 Two circles, each of radius $3$, are drawn with centers at $(0, 16)$, and $(5, 8)$. A line passing through $(3,0)$ is such that the total area of the parts of the three circles to one side of the line is equal to the total area of the parts of the three circles to the other side of it. What is the absolute value of the slope of this line?

 Note that this is equivalent to finding a line such that the distance from $(0,16)$ to the line is the same as the distance from $(5,8)$ to the line. Let the line be $y - ax - b = 0$. Then, we have that:
 \[\frac{|-5a + 8 - b|}{\sqrt{a^2+1}}= \frac{|16 - b|}{\sqrt{a^2+1}} \Longleftrightarrow |-5a+8-b| = |16-b|\]We can split this into two cases.

 Case 1:  $16-b = -5a + 8 - b \Longleftrightarrow a = -\frac{8}{5}$

 In this case, the absolute value of the slope of the line won’t be an integer, and since this is an AIME problem, we know it’s not possible.

 Case 2:  $b-16 = -5a + 8 - b \Longleftrightarrow 2b + 5a = 24$

 But we also know that it passes through the point $(3,0)$, so $-3a-b = 0 \Longleftrightarrow b = -3a$. Plugging this in, we see that $2b + 5a = 24 \Longleftrightarrow a = -24$. $\boxed{24}$.

 



\textbf{Solution 3 (non-rigorous)}

Consider the region of the plane between $x=16$ and $x=17$. The parts of the circles centered at $(14,92)$ and $(19,84)$ in this region have equal area. This is by symmetry- the lines defining the region are 2 units away from the centers of each circle and therefore cut off congruent segments. We will draw the line in a way that uses this symmetry and makes identical cuts on the circles. Since $(17,76)$ is $8$ units below the center of the lower circle, we will have the line exit the region $8$ units above the center of the upper circle, at $(16,100)$. We then find that the slope of the line is $-24$ and our answer is $\boxed{024}$.

 (Note: this solution does not feel rigorous when working through it, but it can be checked easily. In the above diagram, the point $M$ is marked. Rotate the aforementioned region of the plane $180^\circ$ about point $M$, and the fact that certain areas are equal is evident.)

 



\textbf{Solution 4}

We can redefine the coordinate system so that the center of the center circle is the origin, for easier calculations, as the slope of the line and the congruence of the circles do not depend on it. $O_1=(-3, 16)$ $O_2=(0,0)$, and $O_3=(2,8)$. A line bisects a circle iff it passes through the center. Therefore, we can ignore the bottom circle because it contributes an equal area with any line. A line passing through the centroid of any plane system with two perpendicular lines of reflectional symmetry bisects it. We have defined two points of the line, which are $(0,0)$ and $(-\frac{1}{2},12)$. We use the slope formula to calculate the slope, which is $-24$, leading to an answer of $\boxed{024}$. $QED \blacksquare$

 Solution by \href{https://artofproblemsolving.com/wiki/index.php?title=User:A1b2}{a1b2}

 



\textbf{Solution 5 (Easy to understand)}

Notice that any line that passes through the bottom circle's center cuts it in half, so all we really care about are the top two circles. Suppose $\ell$ is the desired line. Draw lines $\ell_1$ and $\ell_2$ both parallel to $\ell$ such that $\ell_1$ passes through $(14,92)$ and $\ell_2$ passes through $(19,84)$. Clearly, $\ell$ must be the "average" of $\ell_1$ and $\ell_2$. Suppose $\ell:=y=mx+b, \ell_1:=y=mx+c, \ell_2:=y=mx+d$. Then $b=76-17m, c=92-14m, d=84-19m$. So we have that 
 \[76-17m=\frac{92-14m+84-19m}{2},\]which yields $m=-24$ for an answer of $\boxed{024}$.

 


	\item (\textit{AIME I 1984}) The function f is defined on the set of integers and satisfies
$f(n)= \begin{cases}  n-3 & \mbox{if }n\ge 1000 \\  f(f(n+5)) & \mbox{if }n<1000 \end{cases}$


 Find $f(84)$.


 
 \textbf{Answer: }997



\textbf{Solution 1}

Define $f^{h} = f(f(\cdots f(f(x))\cdots))$, where the function $f$ is performed $h$ times. We find that $f(84) = f(f(89)) = f^2(89) = f^3(94) = \ldots f^{y}(1004)$. $1004 = 84 + 5(y - 1) \Longrightarrow y = 185$. So we now need to reduce $f^{185}(1004)$.

 Let’s write out a couple more iterations of this function:
 \begin{align*}f^{185}(1004)&=f^{184}(1001)=f^{183}(998)=f^{184}(1003)=f^{183}(1000)\\ &=f^{182}(997)=f^{183}(1002)=f^{182}(999)=f^{183}(1004)\end{align*}
So this function reiterates with a period of 2 for $x$. It might be tempting at first to assume that $f(1004) = 1001$ is the answer; however, that is not true since the solution occurs slightly before that. Start at $f^3(1004)$:
 \[f^{3}(1004)=f^{2}(1001)=f(998)=f^{2}(1003)=f(1000)=\boxed{997}\]

 Note that we should also be suspicious if our answer is $1001$- it is a $4$-digit number, and we were not asked to, say, divide our number by $13$.

 



\textbf{Solution 2}

We start by finding values of the function right under $1000$ since they require iteration of the function.

 
 \[f(999)=f(f(1004))=f(1001)=998\]
 \[f(998)=f(f(1003))=f(1000)=997\]
 \[f(997)=f(f(1002))=f(999)=998\]
 \[f(996)=f(f(1001))=f(998)=997\]

 Soon we realize the $f(k)$ for integers $k<1000$ either equal $998$ or $997$ based on its parity. (If short on time, a guess of $998$ or $997$ can be taken now.)If $k$ is even, $f(k)=997$. If $k$ is odd, $f(k)=998$. $84$ has even parity, so $f(84)=997$.The result may be rigorously shown through induction.

 



\textbf{Solution 3}

Assume that $f(84)$ is to be performed $n+1$ times. Then we have
 \[f(84)=f^{n+1}(84)=f(f^n(84+5))\]In order to find $f(84)$, we want to know the smallest value of
 \[f^n(84+5)\ge1000\]Because then
 \[f(84)=f(f^n(84+5))=(f^n(84+5))-3\]From which we'll get a numerical value for $f(84)$.

 Notice that the value of $n$ we expect to find is basically the smallest $n$ such that after $f(x)=f(f(x+5))$ is performed $\frac{n}{2}$ times and then $f(x)=x-3$ is performed back $\frac{n}{2}$ times, the result is greater than or equal to $1000$. 

 In this case, the value of $n$ for $f(84)$ is $916$, because
 \[84+\frac{916}{2}\cdot5-\frac{916}{2}\cdot3=1000\Longrightarrow f^{916}(84+5))=1000\]Thus
 \[f(84)=f(f^{916}(84+5))=f(1000)=1000-3=\boxed{997}\]

 



\textbf{Solution 4 (really a solution, DO THIS ON A REAL TEST)}

Open up a coding IDE and use \href{https://artofproblemsolving.com/wiki/index.php?title=Recursion}{recursion} to do this problem. The idea is to define a function (I called it $f$, you can call it whatever you want) with parameter $n$ (or 84 in this case) and say if $n$ is greater than 1000, then return $n-3$. Else, return $f(f(n + 5))$. Python code:

 \begin{verbatim}
def f(n):       if n >= 1000:           return n - 3       else:           return f(f(n + 5))   print(f(84))
\end{verbatim}


Or [python]def f(n):

 \begin{verbatim}
if n < 1000:       return f(f(n + 5))  else:      return n-3
\end{verbatim}

print(f(84))[/python]

 Note:"Do THIS ON A REAL TEST" does not mean open up python on the real AIME(hopefully)

 


	\item (\textit{AIME I 1984}) The equation $z^6+z^3+1=0$ has complex roots with argument $\theta$ between $90^\circ$ and $180^\circ$ in the complex plane. Determine the degree measure of $\theta$.


 
 \textbf{Answer: }160



\textbf{Solution 1}

We shall introduce another factor to make the equation easier to solve. If $r$ is a root of $z^6+z^3+1$, then $0=(r^3-1)(r^6+r^3+1)=r^9-1$. The polynomial $x^9-1$ has all of its roots with \href{https://artofproblemsolving.com/wiki/index.php?title=Absolute\_value}{absolute value} $1$ and argument of the form $40m^\circ$ for integer $m$ (the ninth degree \href{https://artofproblemsolving.com/wiki/index.php?title=Roots\_of\_unity}{roots of unity}). Now we simply need to find the root within the desired range that satisfies our original equation $x^6 + x^3 + 1 = 0$.

 This reduces $\theta$ to either $120^{\circ}$ or $160^{\circ}$. But $\theta$ can't be $120^{\circ}$ because if $r=\cos 120^\circ +i\sin 120^\circ$, then $r^6+r^3+1=3$. (When we multiplied by $r^3 - 1$ at the beginning, we introduced some extraneous solutions, and the solution with $120^\circ$ was one of them.) This leaves $\boxed{\theta=160}$.

 \textbf{Note} When evaluation of the LHS we see that the denominator is $z^3 - 1$, which cannot equal 0. Therefore, all third roots of unity are extraneous. One of these roots is $120^\circ$, and I hope that makes it easier to understand. I did not create the above solution btw.

 



\textbf{Solution 2}

The substitution $y=z^3$ simplifies the equation to $y^2+y+1 = 0$. Applying the quadratic formula gives roots $y=-\frac{1}{2}\pm \frac{\sqrt{3}i}{2}$, which have arguments of $120$ and $240,$ respectively. We can write them as $z^3 = \cos 240^\circ + i\sin 240^\circ$ and $z^3 = \cos 120^\circ + i\sin 120^\circ$. So we can use \href{https://artofproblemsolving.com/wiki/index.php?title=De\_Moivre\%27s\_theorem}{De Moivre's theorem} (which I would suggest looking at if you never heard of it before) to find the fractional roots of the expressions above! For $\cos 240^\circ + i\sin 240$ we have $(\cos 240^\circ + i\sin 240^\circ)^{1/3}$ $\Rightarrow$ $\cos 80^\circ + i\sin 80^\circ, \cos 200^\circ + i\sin200^\circ,$ and $\cos 320^\circ + i\sin320^\circ.$Similarly for $(\cos 120^\circ + i\sin 120^\circ)^{1/3}$, we have $\cos 40^\circ + i\sin 40^\circ, \cos 160^\circ + i\sin 160^\circ,$ and $\cos 280^\circ + i\sin 280^\circ.$ The only argument out of all these roots that fits the description is $\theta = \boxed{160}$

 Note: We can add $120$ to the angles of the previous solutions to get new solutions because De Moivre's formula says that $(\cos \theta + i \sin \theta)^n = \cos n\theta + i \sin n\theta$ and $\frac{360}{3} = 120$. ~programmeruser

 



\textbf{Solution 3}

As in Solution 2, make the substitution $u=z^3$. Then we are left with $u^2+u+1=0$, and we would have the solutions $e^{\frac{2\pi}{3}i},e^{\frac{4\pi}{3}i},e^{\frac{8\pi}{3}i},e^{\frac{10}{3}\pi}$. The latter two solutions are obtained by adding one extra revolution around the unit circle. (Notice how we omitted $e^{\frac{6\pi}{3}i}$, since this would yield $e^{2i\pi}=1$ and does not satisfy the equation.) Now, we substitute back, which gives us $z^3=e^{\frac{2\pi}{3}i},e^{\frac{4\pi}{3}i},e^{\frac{10}{3}\pi}\implies z=e^{\frac{2\pi}{9}i},e^{\frac{4\pi}{9}i},e^{\frac{8\pi}{9}i},e^{\frac{10\pi}{9}i}$. The only root in the range $\frac{\pi}{2}<\theta<\pi$ is achieved when $\theta=\frac{8\pi}{9}=\boxed{160^\circ}$.Note that there are 6 solutions to $z$ in this equation, and they are obtained by simply adding more revolutions around the unit circle and dividing by 3.

 


	\item (\textit{AIME I 1984}) In tetrahedron $ABCD$, edge $AB$ has length 3 cm. The area of face $ABC$ is $15\mbox{cm}^2$ and the area of face $ABD$ is $12 \mbox { cm}^2$. These two faces meet each other at a $30^\circ$ angle. Find the volume of the tetrahedron in $\mbox{cm}^3$.


 
 \textbf{Answer: }020



\textbf{Solution 1}


                    \begin{figure}[H]
                        \centering
                        \includegraphics[width=0.5\textwidth]{Solution/AIME_I_1984_9_Solution_Figure_1.png}
                    \end{figure}
                    

 Position face $ABC$ on the bottom. Since $[\triangle ABD] = 12 = \frac{1}{2} \cdot AB \cdot h_{ABD}$, we find that $h_{ABD} = 8$. Because the problem does not specify, we may assume both $ABC$ and $ABD$ to be isosceles triangles. Thus, the height of $ABD$ forms a $30-60-90$ with the height of the tetrahedron. So, $h = \frac{1}{2} (8) = 4$. The volume of the tetrahedron is thus $\frac{1}{3}Bh = \frac{1}{3} \cdot15 \cdot 4 = \boxed{020}$.

 



\textbf{Solution 2 (Rigorous)}

It is clear that $DX=8$ and $CX=10$ where $X$ is the foot of the perpendicular from $D$ and $C$ to side $AB$. Thus $[DXC]=\frac{ab\sin{c}}{2}=20=5 \cdot h \rightarrow h = 4$ where h is the height of the tetrahedron from $D$. Hence, the volume of the tetrahedron is $\frac{bh}{3}=15\cdot \frac{4}{3}=\boxed{020}$~ kilobyte144

 (Note this actually isn't rigorous because they never proved that the height from $D$ to $XC$ is the altitude of the tetrahedron.

 



\textbf{Solution 2.1 (Another Rigorous Solution)}

Drop an altitude from $C$ to $X$ on triangle $ABC$. Through definition of height of a tetrahedron, the height from $D$ down to plane $ABC$ always intersects the altitude of the triangle $ABC$, and the reason this works is due to the shared side and 30 degree measure of the face along the base. Therefore, label that intersection point $Y$. Let $XY$ be $p$. Drop the altitude from $D$ to $AB$ on triangle $ADB$, call it $Z$. Let $XZ = q$. Note that $DZ = 8$

 We show that $XZ = XY$. We know that $XYZ$ is a right angled triangle. then, call angle $XYZ ~\theta$. Then, $\tan(\theta) = q/p$. Also, $\cos(\theta) = \frac{p}{\sqrt{p^2+q^2}}$ and $\sin(\theta) = \frac{q}{\sqrt{p^2+q^2}}$. Squaring we have $\frac{p^2}{p^2+q^2}$ and $\frac{q^2}{p^2+q^2}$. When we equate these we obtain $p^2 = q^2$ and $p = q$, proving the above.

 Therefore, Triangle $DZY$ is right with a 30 degree angle at $Z$. Then, the hypotenuse is 8, and so the smallest side is 4. The volume is then $\frac{1}{3} \cdot 15 \cdot 4 = \boxed{020}$.

 



\textbf{Solution 3 (Sketchy)}

Make faces $ABC$ and $ABD$ right triangles. This makes everything a lot easier. Then do everything in solution 1.

 



\textbf{Solution 4 (coord/vector bash)}

We can use 3D coordinates.

 Let $A = (0, 0, 0)$ and $B = (3, 0, 0).$ WLOG, let $D = \left(\frac{3}{2}, 8, 0\right)$, because the area of $\Delta{ABD} = 12$ and the tetrahedron area won't change if we put it somewhere else with $y=8.$

 To find $C$, we can again let the $x$-coordinate be $\frac{3}{2}$ for simplicity. Note that $C$ is $10$ units away from $AB$ because the area of $\Delta{ABC}$ is $15$. Since the angle between $ABD$ and $ABC$ is $30^\circ$, we can form a 30-60-90 triangle between $A$, $B$, and an altitude dropped from $C$ onto face $ABD$. Since $10$ is the hypotenuse, we get $5\sqrt{3}$ and $5$ as legs. Then $y=5\sqrt{3}$ and $z=5$, so $C = \left(\frac{3}{2}, 5\sqrt{3}, 5\right).$

 (I highly advise you to draw both the tetrahedron and 30-60-90 triangle to get a better perspective.)

 Now, we can move onto vectors. To find the volume of the tetrahedron, we use the formula $\frac{1}{3}Bh.$ Letting $\Delta{ABC}$ be the base we have $B = 15$ (from the problem statement). We need to find the distance between $D$ and $ABC$, and to do this, we should find the projection of $D$ onto face $ABC$.

 Note that we can simplify this to projecting $D$ onto $\mathbf{\overrightarrow{C}}.$ This is because we know the projection will have the same $x$-coordinate as $D$ and $C$, as both are $\frac{3}{2}.$ Now we find $\text{proj}_{\mathbf{\overrightarrow{D}}} \mathbf{\overrightarrow{C}}$, or plugging in our coordinates, $\text{proj}_{\langle\frac{3}{2}, 5\sqrt{3}, 5\rangle} \left\langle\frac{3}{2}, 8, 0\right\rangle$. 

 Let the $x$-coordinates for both be $0$ for simplicity, because we can always add a $\frac{3}{2}$ at the end. Using the projection formula, we get 
 \[\langle 0, 6, 2\sqrt{3}\rangle.\]

 Finally, we calculate the distance between $\left(\frac{3}{2}, 6, 2\sqrt{3}\right)$ and $D$ to be $4$. So the height is $4$, and plugging into our tetrahedron formula we get 
 \[\frac{1}{3}\cdot 15\cdot 4 = \boxed{20}.\]

 -kilobyte144

 


	\item (\textit{AIME I 1984}) Mary told John her score on the American High School Mathematics Examination (AHSME), which was over $80$. From this, John was able to determine the number of problems Mary solved correctly. If Mary's score had been any lower, but still over $80$, John could not have determined this. What was Mary's score? (Recall that the AHSME consists of $30$ multiple choice problems and that one's score, $s$, is computed by the formula $s=30+4c-w$, where $c$ is the number of correct answers and $w$ is the number of wrong answers. Students are not penalized for problems left unanswered.)


 
 \textbf{Answer: }119



\textbf{Solution 1 (Inequalities)}

Let Mary's score, number correct, and number wrong be $s,c,w$ respectively. Then
 \begin{align*} s&=30+4c-w \\ &=30+4(c-1)-(w-4) \\ &=30+4(c+1)-(w+4). \end{align*}
Therefore, Mary could not have left at least five blank; otherwise, one more correct and four more wrong would produce the same score. Similarly, Mary could not have answered at least four wrong (clearly Mary answered at least one right to have a score above $80$, or even $30$.)

 It follows that $c+w\geq 26$ and $w\leq 3$, so $c\geq 23$ and $s=30+4c-w\geq 30+4(23)-3=119$. So Mary scored at least $119$. To see that no result other than $23$ right/$3$ wrong produces $119$, note that $s=119\Rightarrow 4c-w=89$ so $w\equiv 3\pmod{4}$. But if $w=3$, then $c=23$, which was the result given; otherwise $w\geq 7$ and $c\geq 24$, but this implies at least $31$ questions, a contradiction. This makes the minimum score $\boxed{119}$.

 



\textbf{Solution 2 (Arithmetic)}

A less technical approach that still gets the job done:

 Pretend that the question is instead a game, where we are trying to get certain numbers by either adding $4$ or $5.$ The maximum number we can get is $70.$ The goal of the game is to find out what number we can achieve in only ONE method, while all other numbers above that can be achieved with TWO or MORE methods. (Note: This is actually the exact same problem as the original, just reworded differently and now we are adding the score. If this is already confusing, I suggest not looking further.)

 For example, the number $"21"$ can be achieved with only $1$ method $(4+4+4+4+5).$ However, $25$, which is a larger number than $21$, can be achieved with multiple methods (e.g. $5 \cdot 5$ or $4 \cdot 5 + 5$), hence $21$ is not the number we are trying to find. 

 If we make a table of adding $4$ or adding $5$, we will see we get $4, 8, 12, 16, 20,$ etc. if we add only $4$s and if we add $5$ to those numbers then we will get $9, 13, 17, 21, 25,$ etc. Now a key observation to getting this problem correct is that if we can add one of those previous base numbers to $20$, then there will be multiple methods (because $20 = 4 \cdot 5 = 5 \cdot 4$).

 Hence, the number we are looking for cannot be $20$ plus one of those base numbers. Instead, it must be $10$ plus that base number, because that results in the same last digit while maintaining only one method to solve. For example, if we start with $4$, the number $14$ would have only $1$ method to solve, but the number $24$ would have multiple (because $4 + 20 = 24$ and we are trying to avoid adding $20$). The largest number we see that is in our base numbers is $21.$ Hence, our maximum number is $21 + 10 = 31.$

 Note that if we have the number $25$, that can be solved via multiple methods, and if we keep repeating our cycle of base numbers, we are basically adding $20$ to a previous base number, which we don't want.

 And since the maximum number of this game is $31$, that is the number we subtract from the maximum score of $150$, so we get $150 - 31 = \boxed{119}.$

 P.S. didn't think the solution would be this complicated when I first wrote it but it's quite complicated. Look to solution $1$ if you want a concise method using inequalities that's probably better than this solution.

 



\textbf{Solution 3 (Table)}

Based on the value of $c,$ we construct the following table:
 \[\begin{array}{c||c|c|c|c|c|c|c|c|c|c|c|c|c|c|c|c|c|c|c|c}    &\hspace{5.5mm}&\hspace{5.5mm}&\hspace{5.5mm}&\hspace{5.5mm}&\hspace{5.5mm}&\hspace{5.5mm}&\hspace{5.5mm}&&&&&&&&&&&&& \\ [-2.5ex] \boldsymbol{c} &\boldsymbol{\cdots}&\boldsymbol{12}&\boldsymbol{13}&\boldsymbol{14}&\boldsymbol{15}&\boldsymbol{16}&\boldsymbol{17}&\boldsymbol{18}&\boldsymbol{19}&\boldsymbol{20}&\boldsymbol{21}&\boldsymbol{22}&\boldsymbol{23}&\boldsymbol{24}&\boldsymbol{25}&\boldsymbol{26}&\boldsymbol{27}&\boldsymbol{28}&\boldsymbol{29}&\boldsymbol{30} \\  \hline \hline &&&&&&&&&&&&&&&&&&& \\ [-2.25ex] \boldsymbol{s_{\min}} &\cdots&60&65&70&75&80&85&90&95&100&105&110&115&120&125&130&135&140&145&150 \\ \hline   &&&&&&&&&&&&&&&&&&& \\ [-2.25ex] \boldsymbol{s_{\max}} &\cdots&78&82&86&90&94&98&102&106&110&114&118&122&126&130&134&138&142&146&150 \end{array}\]For a fixed value of $c,$ note that $s_{\min}$ occurs at $w=30-c,$ and $s_{\max}$ occurs at $w=0.$ Moreover, all integers from $s_{\min}$ through $s_{\max}$ are attainable. <i>\textbf{To find Mary's score, we look for the lowest score $\boldsymbol{s}$ such that $\boldsymbol{s\geq80}$ and $\boldsymbol{s}$ is contained in exactly one interval.}</i>

 Let $S(c)$ denote the interval of all possible scores $s$ with $c$ correct answers. We need: 

 <ol style="margin-left: 1.5em;"> \item $S(c)\not\subset S(c-1)\cup S(c+1).$ 

 \item $s\in S(c)$ but $s\not\in S(c-1)\cup S(c+1).$

 </ol>It follows that the least such value of $c$ is $23,$ from which the lowest such score $s$ is $\boxed{119}.$

 



\textbf{Solution 4 (Less Rigorous Version of Solution 1)}

Given that Mary's score is $30+4c-w$, two other ways to get that score are $30+4(c+1)-(w+4)$ and $30+4(c-1)-(w-4)$.Since it is clear that $c>1$, we must have $w<4$. In order to minimize the score, assume that $w=3$.The number of problems left blank must be less than $5$ because of the $30+4(c+1)-(w+4)$ case. In order to minimize the score, assume that the number of problems left blank is $4$, making the number of correct problems $23$.Substituting, we get that $s=30+23{\,\cdot\,}4-3$, so $s=\boxed{119}$.

 


	\item (\textit{AIME I 1984}) A gardener plants three maple trees, four oaks, and five birch trees in a row. He plants them in random order, each arrangement being equally likely. Let $\frac m n$ in lowest terms be the probability that no two birch trees are next to one another. Find $m+n$.


 
 \textbf{Answer: }106



\textbf{Solution 1}

First notice that there is no difference between the maple trees and the oak trees; we have only two types, birch trees and "non-birch" trees. (If you don't believe this reasoning, think about it. You could also differentiate the tall oak trees from the short oak trees, and the maple trees with many branches as opposed to those with few branches. Indeed, you could keep dividing until you have them each in their own category, but in the end it will not change the probability of the birch trees being near each other. That is, in the end, you multiply the numerator by the number of ways to arrange the oak and maple trees and you also multiply the denominator by the number of ways to arrange the oak and maple trees, making them cancel out.)

 The five birch trees must be placed amongst the seven previous trees. We can think of these trees as 5 dividers of 8 slots that the birch trees can go in, making ${8\choose5} = 56$ different ways to arrange this.

 There are ${12 \choose 5} = 792$ total ways to arrange the twelve trees, so the probability is $\frac{56}{792} = \frac{7}{99}$.

 The answer is $7 + 99 = \boxed{106}$.

 


 \textbf{(Another way to think about Solution 1)} So you can think of placing the $5$ birch trees and the other trees with the restrictions as described above. Then let's take out one tree between each pair of birch trees. So you would remove $4$ trees that aren't birch. What you are left with is a unique arrangement of $5$ birch trees and $3$ other trees that is unrestricted. Some birch trees might become adjacent after you remove $4$ trees. Adding a tree between each pair of people gives a unique arrangement of $5$ nonadjacent birch trees. This guarantees that there are no adjacent trees. The number of unrestricted $7$ tree arrangments is ${8\choose5} = 56$. Then proceed as told in Solution $1$. 

 



\textbf{Solution 2}

Let $b$, $n$ denote birch tree and not-birch tree, respectively. Notice that we only need $4$ $n$s to separate the $5$ $b$s. Specifically,
 \[b,n,b,n,b,n,b,n,b\]Since we have $7$ $n$s, we are placing the extra $3$ $n$s into the $6$ intervals beside the $b$s.

 Now doing simple casework.

 If all $3$ $n$s are in the same interval, there are $6$ ways.

 If $2$ of the $3$ $n$s are in the same interval, there are $6\cdot5=30$ ways.

 If the $n$s are in $3$ different intervals, there are ${6 \choose 3} =20$ ways.

 In total there are $6+30+20=56$ ways.

 There are ${12\choose5}=792$ ways to distribute the birch trees among all $12$ trees.

 Thus the probability equals $\frac{56}{792}=\frac{7}{99}\Longrightarrow m+n=7+99=\boxed{106}$.

 



\textbf{Solution 3 (using PIE)}

Note that the requested probability is computed by dividing the number of configurations with no adjacent Birch trees by the total number of configurations. We can compute the number of configurations with no adjacent Birch trees using complementary counting and then the Principle of Inclusion-Exclusion.

 The number of configurations with no adjacent Birch trees is equal to the total number of configurations minus the number of configurations with at least one pair of adjacent Birch trees. 

 The total number of configurations is given by $\frac{12!}{3! \cdot 4! \cdot 5!}$. To compute the number of configurations with at least one pair of adjacent Birch trees, we use PIE.

 $\#$(configurations with at least one pair of adjacent Birch trees) $=$ $\#$(configurations with one pair) $-$ $\#$(configurations with two pairs) $+$ $\#$(configurations with three pairs) $-$ $\#$(configurations with four pairs).

 To compute the first term, note that we can treat the adjacent pair of Birch trees as one separate tree. This then gives $\frac{11!}{3! \cdot 3! \cdot 4!}$ configurations. 

 For the second term, we have two cases. The two pairs could either happen consecutively (BBB) or separately (BB BB). They both give $\frac{10!}{2! \cdot 3! \cdot 4!}$ cases. So our second term is $\frac{2 \cdot 10!}{2! \cdot 3! \cdot 4!}$.

 The third term can also happen in two ways. The three pairs could be arranged like BBBB or BBB BB. Both cases together give $\frac{2 \cdot 9!}{3! \cdot 4!}$ arrangements.

 The final term can happen in one way (BBBBB). This gives $\frac{8!}{3! \cdot 4!}$ arrangements.

 Substituting these into our PIE expression, we find that there are $25760$ configurations with at least one pair of adjacent Birch trees. Therefore, there are a total of $\frac{12!}{3! \cdot 4! \cdot 5!} - 25760 = 1960$ configurations with no adjacent Birch trees.

 Thus, the probability of a given configuration having no two adjacent Birch trees is given by $\frac{1960}{\frac{12!}{3! \cdot 4! \cdot 5!}} = \frac{7}{99}$.

 
Therefore, the desired result is given by $7+99 = \boxed{106}$.

 



\textbf{Solution 4}

Here is a solution leaving out nothing. This solution is dedicated to those that are in self study and wish to learn the most they can. I will make it as elementary as possible and intuition based.Arrange first the $3$ maple and $4$ oaks as $MMMOOOO$. We then notice that for none of the $5$ birch trees to be adjacent, they must be put in between these $M$'s and $O$'s. We then see that there are $8$ spots to put these $5$ birch trees in. So we can select $5$ spots for these birch trees in $\binom{8}{5}$. But then, we can rearrange the $M$'s and $O$'s in $7!/(3!4!)=\binom{7}{3}$ ways. So then there are $\binom{8}{5}\binom{7}{3}$ valid arrangements with no given consecutive birch trees.There are then a total of $\frac{12!}{3!4!5!}$ different total arrangements. Therefore the probability is given as $\frac{\binom{8}{5}\binom{7}{3}}{\frac{12!}{3!4!5!}}=\frac{7}{99}$, so the answer is $7+99=\boxed{106}$.

 



\textbf{Solution 5}

Let's remove all of the birch trees so that were left with $3$ maple trees and $4$ oak trees. This can be arranged as $MMMOOOO$. Now, we can place the B's between the letters and at the ends. So, in total we have $8$ places to place the $B$'s and we have $5$ $B$'s to place so in total we have $\dbinom{8}{5}$ ways. Next, we have $\dbinom{12}{5}$ ways to arrange the trees, so our probability is 
 \[\dfrac{\dbinom{8}{5}}{\dbinom{12}{5}} = \dfrac{7}{99} \Longrightarrow \boxed{106}\]

 -jb2015007

 


	\item (\textit{AIME I 1984}) A function $f$ is defined for all real numbers and satisfies $f(2+x)=f(2-x)$ and $f(7+x)=f(7-x)$ for all $x$. If $x=0$ is a root for $f(x)=0$, what is the least number of roots $f(x)=0$ must have in the interval $-1000\leq x \leq 1000$?


 
 \textbf{Answer: }401



\textbf{Solution 1}

If $f(2+x)=f(2-x)$, then substituting $t=2+x$ gives $f(t)=f(4-t)$. Similarly, $f(t)=f(14-t)$. In particular,
 \[f(t)=f(4-t)=f(14-(4-t))=f(t+10)\]

 Since $0$ is a root, all multiples of $10$ are roots, and anything congruent to $4\pmod{10}$ are also roots. To see that these may be the only integer roots, observe that the function
 \[f(x) = \sin \frac{\pi x}{10}\sin \frac{\pi (x-4)}{10}\]satisfies the conditions and has no other roots.

 In the interval $-1000\leq x\leq 1000$, there are $201$ multiples of $10$ and $200$ numbers that are congruent to $4 \pmod{10}$, therefore the minimum number of roots is $\boxed{401}$.

 



\textbf{Solution 2 (non-rigorous)}

We notice that the function has reflectional symmetry across both $x=2$ and $x=7$. We also use the fact that $x=0$ is a root. This shows that $x=4$ and $x=14$ are also roots. We then apply the reflection across the other axis to form $x=\pm 10$ as roots. Continuing this shows that the roots are $0 \mod 10$ or $4 \mod 10$. There are 200 positive roots and 200 negative roots. 0 is also a root, and adding these gives a result of $\boxed{401}$. $QED \blacksquare$

 Solution by \href{https://artofproblemsolving.com/wiki/index.php?title=User:A1b2}{a1b2}

 



\textbf{Solution 3}

Since this is a recursive problem, list out the functions f(2) and f(7) and figure out what is equivalent with them. Then find the x values for the functions that are equal to f(2) and f(7). You will notice that it starts at x=0, then it goes to x=5, x=10, etc... each f() has two possible x values, but we are only counting the total number of x values so ignore the double counted values. This means that $x = 0, \pm 5, \pm 10, \pm 15... \pm 1000$ so the answer is 400 + 1 = $\boxed{401}$

 



\textbf{Solution 4}

Let $z$ be an arbitrary zero. If $z=2-x$, then $x=2-z$ and $2+x=4-z$. Repeat with other equation to find if $z$ is a zero then so are $4-z$ and $14-z$. From $0$, we get $4$ and $14$. Now note that applying either of these twice will return $z$, so we must apply them in an alternating fashion for distinct roots. Doing so to $4$ and $14$ returns $10$ and $-10$, respectively. A pattern will emerge of each path hitting a multiple of $10$ after $2$ moves. Hence, we will reach $\pm 1000$ after $200$ jumps in either direction. Including zero, there are $2\cdot200+1=\boxed{401}$

 


	\item (\textit{AIME I 1984}) Find the value of $10\cot(\cot^{-1}3+\cot^{-1}7+\cot^{-1}13+\cot^{-1}21).$


 
 \textbf{Answer: }015



\textbf{Solution 1}

We know that $\tan(\arctan(x)) = x$ so we can repeatedly apply the addition formula, $\tan(x+y) = \frac{\tan(x)+\tan(y)}{1-\tan(x)\tan(y)}$. Let $a = \cot^{-1}(3)$, $b=\cot^{-1}(7)$, $c=\cot^{-1}(13)$, and $d=\cot^{-1}(21)$. We have

 $\tan(a)=\frac{1}{3},\quad\tan(b)=\frac{1}{7},\quad\tan(c)=\frac{1}{13},\quad\tan(d)=\frac{1}{21}$,

 

 so

 $\tan(a+b) = \frac{\frac{1}{3}+\frac{1}{7}}{1-\frac{1}{21}} = \frac{1}{2}$

 

 and

 $\tan(c+d) = \frac{\frac{1}{13}+\frac{1}{21}}{1-\frac{1}{273}} = \frac{1}{8}$,

 

 so

 $\tan((a+b)+(c+d)) = \frac{\frac{1}{2}+\frac{1}{8}}{1-\frac{1}{16}} = \frac{2}{3}$.

 

 Thus our answer is $10\cdot\frac{3}{2}=\boxed{015}$.

 



\textbf{Solution 2}

Apply the formula $\cot^{-1}x + \cot^{-1} y = \cot^{-1}\left(\frac {xy-1}{x+y}\right)$ repeatedly. Using it twice on the inside, the desired sum becomes $\cot (\cot^{-1}2+\cot^{-1}8)$. This sum can then be tackled by taking the cotangent of both sides of the inverse cotangent addition formula shown at the beginning.

 



\textbf{Solution 3}

On the coordinate plane, let $O=(0,0)$, $A_1=(3,0)$, $A_2=(3,1)$, $B_1=(21,7)$, $B_2=(20,10)$, $C_1=(260,130)$, $C_2=(250,150)$, $D_1=(5250,3150)$, $D_2=(5100,3400)$, and $H=(5100,0)$. We see that $\cot^{-1}(\angle A_2OA_1)=3$, $\cot^{-1}(\angle B_2OB_1)=7$, $\cot^{-1}(\angle C_2OC_1)=13$, and $\cot^{-1}(\angle D_2OD_1)=21$. The sum of these four angles forms the angle of triangle $OD_2H$, which has a cotangent of $\frac{5100}{3400}=\frac{3}{2}$, which must mean that $\cot( \cot^{-1}3+\cot^{-1}7+\cot^{-1}13+\cot^{-1}21)=\frac{3}{2}$. So the answer is $10\cdot\left(\frac{3}{2}\right)=\boxed{015}.$

 



\textbf{Solution 4}

Recall that $\cot^{-1}\theta = \frac{\pi}{2} - \tan^{-1}\theta$ and that $\arg(a + bi) = \tan^{-1}\frac{b}{a}$. Then letting $w = 1 + 3i, x = 1 + 7i, y = 1 + 13i,$ and $z = 1 + 21i$, we are left with

 
 \[10\cot(\frac{\pi}{2} - \arg w + \frac{\pi}{2} - \arg x + \frac{\pi}{2} - \arg y + \frac{\pi}{2} - \arg z) = 10\cot(2\pi - \arg wxyz)\]
 \[= -10\cot(\arg wxyz).\]

 Expanding $wxyz$, we are left with
 \[(3+i)(7+i)(13+i)(21+i) = (20+10i)(13+i)(21+i)\]
 \[= (2+i)(13+i)(21+i)\]
 \[= (25+15i)(21+i)\]
 \[= (5+3i)(21+i)\]
 \[= (102+68i)\]
 \[= (3+2i)\]
 \[= 10\cot \tan^{-1}\frac{2}{3}\]
 \[= 10 \cdot \frac{3}{2} = \boxed{015}\]

 


	\item (\textit{AIME I 1984}) What is the largest even integer that cannot be written as the sum of two odd composite numbers?


 
 \textbf{Answer: }038



\textbf{Solution 1}

Take an even positive integer $x$. $x$ is either $0 \bmod{6}$, $2 \bmod{6}$, or $4 \bmod{6}$. Notice that the numbers $9$, $15$, $21$, ... , and in general $9 + 6n$ for nonnegative $n$ are odd composites. We now have 3 cases:

 If $x \ge 18$ and is $0 \bmod{6}$, $x$ can be expressed as $9 + (9+6n)$ for some nonnegative $n$. Note that $9$ and $9+6n$ are both odd composites. 

 If $x\ge 44$ and is $2 \bmod{6}$, $x$ can be expressed as $35 + (9+6n)$ for some nonnegative $n$. Note that $35$ and $9+6n$ are both odd composites. 

 If $x\ge 34$ and is $4 \bmod{6}$, $x$ can be expressed as $25 + (9+6n)$ for some nonnegative $n$. Note that $25$ and $9+6n$ are both odd composites.

 
Clearly, if $x \ge 44$, it can be expressed as a sum of 2 odd composites. However, if $x = 42$, it can also be expressed using case 1, and if $x = 40$, using case 3. $38$ is the largest even integer that our cases do not cover. If we examine the possible ways of splitting $38$ into two addends, we see that no pair of odd composites add to $38$. Therefore, $\boxed{038}$ is the largest possible number that is not expressible as the sum of two odd composite numbers.

 



\textbf{Solution 2}

Let $n$ be an integer that cannot be written as the sum of two odd composite numbers. If $n>33$, then $n-9,n-15,n-21,n-25,n-27,$ and $n-33$ must all be prime (or $n-33=1$, which yields $n=34=9+25$ which does not work). Thus $n-9,n-15,n-21,n-27,$ and $n-33$ form a prime quintuplet. However, only one prime quintuplet exists as exactly one of those 5 numbers must be divisible by 5.This prime quintuplet is $5,11,17,23,$ and $29$, yielding a maximal answer of 38. Since $38-25=13$, which is prime, the answer is $\boxed{038}$.

 



\textbf{Solution 3 (bash)}

Let $2n$ be an even integer. Using the \href{https://artofproblemsolving.com/wiki/index.php?title=Chicken\_McNugget\_Theorem}{Chicken McNugget Theorem} on the two smallest odd composite integers that are relatively prime from each other, 9 and 25, show that the maximum is 191, and the maximum even integer is 190 or lower. We use the fact that sufficiently high multiples of 6, 10, 14, 22, etc. can be represented as $n+n$. We bash each case until we find one that works.

 



\textbf{Solution 4}

Claim: The answer is $\boxed{038}$.

 Proof:It is fairly easy to show 38 can't be split into 2 odd composites.

 Assume there exists an even integer $m > 38$ that m can't be split into 2 odd composites.

 Then, we can consider m modulo 5.

 If m = 0 mod 5, we can express m = 15 + 5k for some integer k. Since $m > 20$, $k > 1$ so $k \geq 2$. Thus, 5k is composite. Since 15 is odd and composite, m is even, 5k should be odd as well. 

 If m = 1 mod 5, we can express m = 21 + 5k for some integer k. Since $m > 26$, $k > 1$ so $k \geq 2$. Thus, 5k is composite. Since 21 is odd and composite, m is even, 5k should be odd as well. 

 If m = 2 mod 5, we can express m = 27 + 5k for some integer k. Since $m > 32$, $k > 1$ so $k \geq 2$. Thus, 5k is composite. Since 27 is odd and composite, m is even, 5k should be odd as well. 

 If m = 3 mod 5, we can express m = 33 + 5k for some integer k. Since $m > 38$, $k > 1$ so $k \geq 2$. Thus, 5k is composite. Since 33 is odd and composite, m is even, 5k should be odd as well. 

 If m = 4 mod 5, we can express m = 9 + 5k for some integer k. Since $m > 14$, $k > 1$ so $k \geq 2$. Thus, 5k is composite. Since 9 is odd and composite, m is even, 5k should be odd as well. 

 Thus, in all cases we can split m into 2 odd composites, and we get at a contradiction. $\blacksquare$

 -Alexlikemath

 



\textbf{Solution 5}

All numbers that could possibly work must be $2 \cdot p$ where $p$ is prime. As previous solutions stated, the maximum number that could possibly work by Chicken McNugget is $9 \cdot 25 - 9 - 25 = 225-34 = 191$. We then bash from top to bottom:

 1. $178 = 89 \cdot 2 => 87 + 91$ - refuted

 2. $166 = 83 \cdot 2 => 81 + 85$ - refuted

 3. $158 = 79 \cdot 2 => 77 + 81$ - refuted

 4. $146 = 73 \cdot 2 => 69 + 77$ - refuted

 5. $142 = 71 \cdot 2 => 65 + 77$ - refuted

 6. $134 = 67 \cdot 2 => 65 + 69$ - refuted

 7. $122 = 61 \cdot 2 => 57 + 65$ - refuted

 8. $118 = 59 \cdot 2 => 55 + 63$ - refuted

 9. $106 = 53 \cdot 2 => 51 + 55$ - refuted

 10. $94 = 47 \cdot 2 => 45 + 49$ - refuted

 11. $86 = 43 \cdot 2 => 35 + 51$ - refuted

 12. $82 = 41 \cdot 2 => 33 + 49$ - refuted

 13. $74 = 37 \cdot 2 => 35 + 39$ - refuted

 14. $62 = 31 \cdot 2 => 27 + 35$ - refuted

 15. $58 = 29 \cdot 2 => 25 + 33$ - refuted

 16. $46 = 23 \cdot 2 => 21 + 25$ - refuted

 17. $38 = 19 \cdot 2 => = 19 + 19$ - it works!

 Because we did a very systematic bash as shown, we are confident the answer is $\boxed {038}$

 



\textbf{Solution 6 (casework)}

As stated above, all numbers that could possibly work must be $2 \cdot p$ where $p$ is prime. If $p$ > 30, we consider $p$ by modulo 30. $p$ could be 1,7,11,13,17,19,23,29 modulo 30. $2 \cdot p$ can be expressed as ($p$+$q$)+($p$-$q$) for some positive, even $q$ less then $p$.

 If $p$ = $1 \bmod{30}$, p±4 would both be composite

 If $p$ = $7 \bmod{30}$, p±2 would both be composite

 If $p$ = $11 \bmod{30}$, p±14 would both be composite

 If $p$ = $13 \bmod{30}$, p±8 would both be composite

 If $p$ = $17 \bmod{30}$, p±8 would both be composite

 If $p$ = $19 \bmod{30}$, p±14 would both be composite

 If $p$ = $23 \bmod{30}$, p±2 would both be composite

 If $p$ = $29 \bmod{30}$, p±4 would both be composite

 So $p$ < 30

 From here, just try all possible p and find the answer is $\boxed {038}$

 


	\item (\textit{AIME I 1984}) Determine $w^2+x^2+y^2+z^2$ if


 $\frac{x^2}{2^2-1}+\frac{y^2}{2^2-3^2}+\frac{z^2}{2^2-5^2}+\frac{w^2}{2^2-7^2}=1$

 $\frac{x^2}{4^2-1}+\frac{y^2}{4^2-3^2}+\frac{z^2}{4^2-5^2}+\frac{w^2}{4^2-7^2}=1$

 $\frac{x^2}{6^2-1}+\frac{y^2}{6^2-3^2}+\frac{z^2}{6^2-5^2}+\frac{w^2}{6^2-7^2}=1$

 $\frac{x^2}{8^2-1}+\frac{y^2}{8^2-3^2}+\frac{z^2}{8^2-5^2}+\frac{w^2}{8^2-7^2}=1$

 
 \textbf{Answer: }036



\textbf{Solution 1}

Rewrite the system of equations as 
 \[\frac{x^{2}}{t-1}+\frac{y^{2}}{t-3^{2}}+\frac{z^{2}}{t-5^{2}}+\frac{w^{2}}{t-7^{2}}=1.\] This equation is satisfied when $t \in \{4, 16, 36, 64\}$. After clearing fractions, for each of the values $t=4,16,36,64$, we have the equation 
 \[x^2P_1(t)+y^2P_3(t)+z^2P_5(t)+w^2P_7(t)=F(t),\]where $F(t)=(t-1^2)(t-3^2)(t-5^2)(t-7^2)$ and $P_k(t)=F(t)/(t-k^2)$, for $k=1,3,5,7$.

 Since the polynomials on each side are equal at $t=4,16,36,64$, we can express the difference of the two polynomials by a quartic polynomial that has roots at $t=4,16,36,64$, so
 \begin{align} \tag{\dag}x^2P_1(t)+y^2P_3(t)+z^2P_5(t)+w^2P_7(t)-F(t) = -(t-4)(t-16)(t-36)(t-64) \end{align}The leading coefficient of the RHS is $-1$ because the leading coefficient of the LHS is $-1$.

 Plug in $t=1^2, 3^2, 5^2, 7^2$ in succession, into $(\dag)$. In each case, most terms drop, and we end up with
 \begin{align*}  x^2=\frac{3^2\cdot 5^2\cdot 7^2}{2^{10}}, \quad y^2=\frac{3^3\cdot 5\cdot 7\cdot 11}{2^{10}},\quad z^2=\frac{3^2\cdot 7\cdot 11\cdot 13}{2^{10}},\quad w^2=\frac{3^2\cdot 5\cdot 11\cdot 13}{2^{10}} \end{align*}
Adding them up we get the sum as $3^2\cdot 4=\boxed{036}$.

 \textbf{Postscript for the puzzled}: This solution which is seemingly unnecessarily redundant in that it computes $x^2,y^2,z^2,$ and $w^2$ separately before adding them to obtain the final answer is appealing because it gives the individual values of $x^2,y^2,z^2,$ and $w^2$ which can be plugged into the given equations to check.Rewrite By MESITEMATH

 



\textbf{Solution 2}

As in Solution 1, we have 
 \[x^2P_1(t)+y^2P_3(t)+z^2P_5(t)+w^2P_7(t)=F(t),\]where $F(t)=(t-1^2)(t-3^2)(t-5^2)(t-7^2)$ and $P_k(t)=F(t)/(t-k^2)$, for $k=1,3,5,7$.

 Now the coefficient of $t^3$ on both sides must be equal. So instead of expanding it fully, we will find what the coefficients of the $t^4$ and $t^3$ terms are, so we can eventually apply Vieta's. We can write the long equation as 
 \[(x^2 + y^2 + z^2 + w^2)t^3 + \dots = t^4 - (1^2 + 3^2 + 5^2 + 7^2)t^3 + \dots\] Rearranging gives us 
 \[t^4 - (1^2 + 3^2 + 5^2 + 7^2 + x^2 + y^2 + z^2 + w^2)t^3 \dots = 0.\] By Vieta's, we know that the sum of the roots of this equation is 
 \[1^2 + 3^2 + 5^2 + 7^2 + x^2 + y^2 + z^2 + w^2 = 2^2 + 4^2 + 6^2 + 8^2.\] (recall that the roots of the original and this manipulated form of it had roots $2^2, 4^2, 6^2,$ and $8^2$). Thus, 
 \[x^2 + y^2 + z^2 + w^2 = 2^2 + 4^2 + 6^2 + 8^2 - 1^2 - 3^2 - 5^2 - 7^2 = \boxed{36}.\]

 



\textbf{Solution 3 (Highly Unrecommended)}

Before starting this solution, I highly recommend never following this unless you have no idea what to do with an hour of your time. Even so, learning the above solutions will be more beneficial.

 
 \begin{align*} \frac{x^2}{2^2-1}+\frac{y^2}{2^2-3^2}+\frac{z^2}{2^2-5^2}+\frac{w^2}{2^2-7^2}=1\\ \frac{x^2}{4^2-1}+\frac{y^2}{4^2-3^2}+\frac{z^2}{4^2-5^2}+\frac{w^2}{4^2-7^2}=1\\ \frac{x^2}{6^2-1}+\frac{y^2}{6^2-3^2}+\frac{z^2}{6^2-5^2}+\frac{w^2}{6^2-7^2}=1\\ \frac{x^2}{8^2-1}+\frac{y^2}{8^2-3^2}+\frac{z^2}{8^2-5^2}+\frac{w^2}{8^2-7^2}=1\\ \end{align*}
can be rewritten as
 \begin{align*} \frac{x^2}{3}-\frac{y^2}{5}-\frac{z^2}{21}-\frac{w^2}{45}=1\\ \frac{x^2}{15}+\frac{y^2}{7}-\frac{z^2}{9}-\frac{w^2}{33}=1\\ \frac{x^2}{35}+\frac{y^2}{27}+\frac{z^2}{11}-\frac{w^2}{13}=1\\ \frac{x^2}{63}+\frac{y^2}{55}+\frac{z^2}{39}+\frac{w^2}{15}=1\\ \end{align*}
You might be able to see where this is going. First off, find $\text{lcm}(3,5,21,45),\text{lcm}(15,7,9,33), \text{lcm}(35,27,11,13),$ and $\text{lcm}(63,55,39,15)$. Then, multiply by the respective lcm to clear all of the denominators. Once you do that, manipulate the equations to solve for $w^2+x^2+y^2+z^2$.

 Now, most of this is just a brainless bash, and reemphasizing, please try to learn the above solutions. This is only a last resort and only to be used if you have too much time left. The exact amount of time this bash takes depends on the person and how quickly they can manipulate the equations.

 



\textbf{Solution 4 (Fast, Efficient, Smart, Highly Recommended, Best)}

Notice how on each line, we have equations of the form $\frac{x^2}{a-1^2}+\frac{y^2}{a-3^2}+\frac{z^2}{a-5^2}+\frac{w^2}{a-7^2}=1$. We can let this be a polynomial, with respect to $a$. We can say that $w^2$, $x^2$, $y^2$, and $z^2$ are coefficients, and not variables. So, we can now expand the fractions to get$(a-1)(a-9)(a-25)(a-49)=x^2(a-9)(a-25)(a-49)$$+ y^2(a-1)(a-25)(a-49)$$+ z^2(a-1)(a-9)(a-49)$$+ w^2(a-1)(a-9)(a-25)$.

 Now, we have arrived at this huge expression, but what do we do with it?

 Well, we can look at what we want to find - $x^2+y^2+z^2+w^2$. So, we want the sum of $x^2$, $y^2$, $z^2$, and $w^2$. Looking back to our expression, we can note how on the right hand side, the $a^3$ terms add to $x^2+y^2+z^2+w^2$. Also, on the left hand side, the $a^3$ coefficient is $-84$ (which is achievable by Vieta's formulas rather than expanding if you want to save a few seconds). So, moving all the $a^3$ terms to the left hand side, then we have that by Vieta's formulas, the sum of the roots is $-84-x^2-y^2-z^2-w^2=-(2^2+4^2+6^2+8^2)$. Then, we can solve to find that $x^2+y^2+z^2+w^2=120-84=\boxed{036}$.

 


\end{enumerate}